% 2011 年北京卷理科第 7 题：四面体的三视图。
% 来源 PDF：普通高考/2011/2011北京理.pdf，SHA-256 79351b2b5d5afd29a2d46d50526284443edad4deb64fbaf841db31ce2570c311。
% 原页：第 1 页（A3 三栏）第 2 栏上部，书页码 647。
% 重绘依据：照原页——正（主）视图为直角边 4、4 的直角三角形（直角在左下）；
%   侧（左）视图为直角边 4（竖直，与正视图等高）、3 的直角三角形（直角在右下）；
%   俯视图为直角边 4、3 的直角三角形（直角在右下）。不按比例。
\begin{tikzpicture}[>=stealth, semithick, scale=0.95,
  dim/.style={<->, thin}]

  % 正（主）视图：直角在左下
  \draw (0,0) -- (0,4) -- (4,0) -- cycle;
  % 侧（左）视图：直角在右下
  \draw (9.4,0) -- (9.4,4) -- (6.4,0) -- cycle;
  % 俯视图：直角在右下
  \draw (0,-5.7) -- (4,-5.7) -- (4,-2.7) -- cycle;

  % 尺寸标注：正视图高 4
  \draw[dim] (-0.75,0) -- (-0.75,4);
  \node[left] at (-0.85,2) {$4$};
  % 尺寸标注：正视图宽 4
  \draw[dim] (0,-0.75) -- (4,-0.75);
  \node[below] at (2,-0.75) {$4$};
  % 尺寸标注：侧视图宽 3
  \draw[dim] (6.4,-0.75) -- (9.4,-0.75);
  \node[below] at (7.9,-0.75) {$3$};

  % 视图名称
  \node at (2,-1.55) {正（主）视图};
  \node at (7.9,-1.55) {侧（左）视图};
  \node at (2,-6.5) {俯视图};
\end{tikzpicture}
